% Generated draft: requires mathematical and source review.
\documentclass[11pt]{article}
\usepackage[margin=25mm]{geometry}
\usepackage{fontspec}
\setmainfont{Latin Modern Roman}
\usepackage{amsmath,amssymb,mathtools,unicode-math}
\setmathfont{Latin Modern Math}
\usepackage{enumitem,needspace,xurl,hyperref}
\setlength{\parindent}{0pt}
\setlength{\parskip}{5pt}
\setlength{\emergencystretch}{3em}
\newcommand{\N}{\mathbb{N}}
\newcommand{\Z}{\mathbb{Z}}
\newcommand{\Q}{\mathbb{Q}}
\newcommand{\R}{\mathbb{R}}
\newcommand{\C}{\mathbb{C}}
\newcommand{\F}{\mathbb{F}}
\newcommand{\A}{\mathbb{A}}
\newcommand{\PP}{\mathbb{P}}
\newcommand{\E}{\mathbb{E}}
\newcommand{\eps}{\varepsilon}
\newcommand{\dd}{\,\mathrm{d}}
\begin{document}
\section*{UnsolvedMath ID: 4200007}
\textbf{Source problem code:} catalogue entry
\subsection{Definitions and mathematical statement}
A planar Birkhoff billiard moves along straight segments inside a table and reflects at a smooth boundary point with equal incidence and reflection angles. A nonsingular orbit avoids corner collisions and tangencies.

(a) Let \(D\subset\mathbb R^2\) be a bounded strictly convex table with smooth boundary. In the usual phase-space formulation let \(M=\partial D\times(0,\pi)\) record a boundary point and a nongrazing outgoing angle, and let \(B:M\to M\) be the billiard map. For each integer \(n\ge1\), put
\[
 \operatorname{Per}_n(B)=\{z\in M:B^n(z)=z\}.
\]
Is \(\operatorname{Per}_n(B)\) nowhere dense in \(M\), that is,
\[
 \operatorname{int}_M\bigl(\overline{\operatorname{Per}_n(B)}^{\,M}\bigr)=\varnothing,
\]
for every such table and every \(n\ge1\)? This uses the period-dividing-\(n\) convention for \(n\)-periodic; the excerpt does not specify least period. The periodic set is interpreted as initial conditions in phase space, since no topology on a quotient space of geometric orbits is supplied. ``Smooth'' is read conventionally as \(C^\infty\); strict convexity is not strengthened to everywhere positive curvature.

(b) Does every polygonal billiard table have at least one nonsingular periodic trajectory? Such a trajectory has positive length, finitely many reflections in side interiors, and returns to the same directed billiard state. The question is not restricted to convex polygons. The conventional setting is a bounded simple planar polygon with nonempty interior; the excerpt does not specify whether multiply connected or other generalized polygonal tables are included.

The supplied summary says (b) is settled for rational polygons, meaning polygons whose interior angles are rational multiples of \(\pi\). It reports an unverified 2026 preprint claiming all polygons, while saying (a) appears unresolved. No authors, title, or reference for that preprint are provided, so its claim is not treated as an established theorem. The original list is from 2000 and no later authoritative resolution check is recorded.
\subsection{Short English statement}
Are periodic initial conditions of each fixed period nowhere dense in every smooth strictly convex planar billiard? Separately, must every polygonal table admit a nonsingular periodic trajectory? The record reports the rational-polygon case settled and a general 2026 proof claim only as an unverified preprint.
\subsection{Sources}
\begin{itemize}
\item Local registry: \url{data/unsolvedmath/problems.json}, permanent UnsolvedMath ID \(4200007\), source code catalogue entry; supplied statement, background, and research summary, updated 2026-08-06.
\item Knill, \emph{Open problems in Hamiltonian dynamics} (2000), item 7, ``Periodic orbits'': \url{https://people.math.harvard.edu/~knill/seminars/intr/}. HTML extraction accessed 2026-07-29; no later authoritative resolution check is recorded.
\item The summary reports the rational-polygon result and an unverified 2026 preprint claiming the general polygon case, without supplying identifiable bibliographic references.
\end{itemize}
\end{document}
